\begin{abstract}
Language models let attackers rewrite phishing cheaply. A gateway detector trained on a fixed partition meets each rewrite as an unannounced shift of the message distribution, while verified labels for the new campaign arrive late and cost analyst time. We ask whether a transformer detector can recognise such shifts without labels and recover its accuracy from a small purchased label budget. SLIDE represents each e-mail by a fine-tuned DistilRoBERTa embedding and four language-model surprisal statistics. Its two-gate drift test applies Holm-corrected Kolmogorov--Smirnov tests to 64 random projections plus the detector score and requires a sliced-Wasserstein effect size above a reference-calibrated threshold; an alarm buys 500 labels, and the head is then retrained with class-balanced replay under representation-space adversarial training. Evaluation is prequential on a 28{,}000-message public stream that moves from seven human corpora to two LLM corpora and a T5 paraphrase campaign, over ten realisations and against nine baselines. SLIDE reached a stream MCC of 0.898 with 2{,}000 labels. Fully supervised ADWIN retraining of the same head also reached 0.898, and online TF-IDF learning 0.905, each with 27{,}500 labels. All 40 segment changes raised an alarm within one batch, with no false alarm. The same 2{,}000 labels bought at random times gave 0.786, whereas an oracle that knows the boundaries gave 0.902. A 400-character benign prefix cut the recall of every transformer detector to at most 0.11. Sliced label-free testing thus times label purchases close to an oracle at 7.3\% of the label cost.
\end{abstract}

\begin{keyword}
Phishing detection \sep Large language models \sep Concept drift \sep Sliced Wasserstein distance \sep Adversarial robustness \sep Data streams
\end{keyword}
